Large language models are costly to serve because their weights and their key--value (KV) cache must fit in GPU memory and because token generation is limited by memory bandwidth. Post-training compression reduces this cost, but published methods are usually applied with one setting for every layer and are judged by simulated perplexity rather than by measurements in the engine that serves the model. This thesis presents a sensitivity-driven framework that addresses both points. Stage~0 measures each decoder layer by the rise in perplexity when the layer is bypassed and writes per-layer plans; Stages~1 to 3 apply published quantization, pruning and KV-cache techniques, each on a uniform plan (Standard) and on the Stage~0 plan (Our method) with the same code; Stage~4 measures the model in Hugging Face Transformers, llama.cpp and vLLM; and a constrained Bayesian search chooses the techniques and settings under limits on perplexity and decode latency. On TinyLlama (1.1 billion parameters), sensitivity is concentrated in five of the 22 layers. Giving 8 bits to the sensitive layers lowers the perplexity of all eleven weight-quantization techniques compared with 4 bits in every layer, for a file of 1.05~GB instead of 0.81~GB, and leaving them unpruned lowers the perplexity of all three pruning techniques at every level. The plan does not help where it spends fewer bits than the uniform setting or, in pruning, at equal size. Pruned models show no gain in size or latency in the engines, and models with rounded activations are 13 to 17 times slower than the uncompressed model in vLLM. The search produced two final models that are about 60 percent smaller and 20 percent faster than the uncompressed model in vLLM: OmniQuant weight-only quantization with light unstructured pruning and InfiniGen (0.838~GB, 1.98~ms per token), and round-to-nearest weight-only quantization with SnapKV (0.898~GB, 2.00~ms per token), against 2.200~GB and 2.54~ms per token. Their mean accuracy on four tasks is 11 and 21 percent lower, which shows that a perplexity limit alone does not protect accuracy.
